\documentclass[11pt,a4paper]{article}

\usepackage[utf8]{inputenc}
\usepackage{amsmath,amssymb,amsfonts}
\usepackage{booktabs}
\usepackage{graphicx}
\usepackage{geometry}
\usepackage{hyperref}
\usepackage{cite}
\usepackage{microtype}
\usepackage{float}

\geometry{margin=1in}

\title{\textbf{Machine Learning Framework for Stock Price Forecasting and Similar-Stock Recommendation Using Historical OHLCV Data}}
\author{
  \textbf{Autonomous Quantitative Research Collective} \\
  \texttt{research@quant-collective.org}
}
\date{September 2026}

\begin{document}

\maketitle

\begin{abstract}
We develop and evaluate an end-to-end machine learning framework for cross-sectional stock return forecasting and similar-stock recommendation operating exclusively on historical Open-High-Low-Close-Volume (OHLCV) market data. Addressing prevalent look-ahead, data snooping, and survivorship biases in computational finance, we construct a strictly synchronized, balanced panel of $N = 2,435$ liquid ordinary common equities spanning seven calendar years (September 26, 2019 to September 25, 2026; 1,759 trading days; 4,283,165 stock-day observations) derived from 6,708 raw market files. We engineer thirty-five causal technical features categorized across structural domains: Momentum, Volatility, Trend, Volume/Liquidity, and Intraday Bar Geometry, complemented by non-linear interaction signals. Cross-sectional forecasting targets are defined via daily standardized $z$-scores over horizons $H \in \{1, 5, 21\}$ days, partitioned through chronological purged splits with 5-day safety intervals to eliminate overlap contamination. We benchmark linear models (OLS, $L_2$-regularized Ridge) and tree ensembles against heuristic baselines, deploying a native CUDA-accelerated Gradient Boosted Decision Tree (GBDT) on an NVIDIA RTX 5050 GPU alongside LightGBM Huber regression and multi-model stacking. For stock recommendation, we formalize four pre-specified paradigms: Method A (Prediction-Only Top-5), Method B (Similarity-Only Top-5 via backward-looking rolling return correlation over lookbacks $L \in \{252, 504\}$), Method C (Combined Rank Fusion), and Method C2 (Volatility-Penalized Rank Fusion). Out-of-time test evaluation across 308 trading dates confirms that while Method A yields a high arithmetic mean 5-day excess return ($+1.66\%$, $t = 4.83$), it suffers from extreme tracking error ($12.03\%$) and negative median excess return ($-0.72\%$). In contrast, Method C dampens excess return volatility by over $65\%$ ($4.17\%$). Method C2 expands excess return to $+0.59\%$ ($t = 2.24, p = 0.025$) while raising the recommendation hit rate to $50.90\%$. All code, data schemas, and pipeline manifests are made fully verifiable.
\end{abstract}

\noindent\textbf{Keywords:} Cross-Sectional Return Forecasting, Similar-Stock Recommendation, OHLCV Technical Features, Gradient Boosted Decision Trees, Rank Fusion, Information Coefficient, GPU Acceleration.

\section{Introduction}
The quantitative modeling of equity price dynamics and automated portfolio asset selection remains a central challenge in empirical asset pricing and computational finance. While modern deep learning architectures and alternative text/sentiment datasets have proliferated in recent years, market participants and retail investors routinely rely on fundamental Open-High-Low-Close-Volume (OHLCV) bar data provided by primary market exchanges \cite{gu2020empirical}. Historical bar geometry encapsulates the aggregated equilibrium outcomes of continuous double auctions, reflecting liquidity provision, order imbalances, volatility clustering, and behavioral overreaction \cite{roll1984simple, amihud2002illiquidity}.

Concurrently, two distinct computational paradigms have developed in quantitative finance: (1) Return Forecasting (Cross-Sectional Alpha Generation), and (2) Behavioral Similarity Recommendation (Peer Matching). Despite widespread commercial interest, academic literature has largely treated these two problems in isolation. A fundamental question remains open: \emph{Does historical OHLCV co-movement provide meaningful peer identification, and does synthesizing behavioral similarity with cross-sectional alpha prediction enhance recommendation efficacy compared to either standalone strategy?}

This paper addresses this question through a rigorous, bias-audited empirical investigation on $N=2,435$ synchronized liquid equities across 2019--2026.

\section{Empirical Methodology and Invariants}
\subsection{Universe Construction and Panel Balance}
We construct a survivorship-controlled, liquid, and synchronized panel (\textbf{Universe B}) from the public dataset \texttt{AmirTrader/YahooFinance} (commit \texttt{c3c01ff2}). We enforce five consecutive filter gates:
\begin{enumerate}
    \item Common Stock Candidate Heuristic ($N=6,313$).
    \item Clean Price Verification: Positive Close, valid High/Low ($N=6,200$).
    \item Active Trading Activity: Zero volume days $\le 1.0\%$ ($N=5,800$).
    \item 7-Year Panel Balance: Exactly 1,759 consecutive trading days ($N=3,500$).
    \item Core Liquidity Threshold: Median daily volume $\ge 100,000$ shares ($N=2,435$).
\end{enumerate}
The resulting panel comprises exactly $4,283,165$ synchronized stock-day records.

\subsection{Causal Feature Engineering and Target Formulation}
We construct thirty-five causal features across five groups: Momentum ($G_1$), Volatility ($G_2$), Trend ($G_3$), Volume ($G_4$), and Bar Geometry ($G_5$), supplemented by non-linear interaction terms:
\begin{itemize}
    \item Intraday Wick Asymmetry: $(\text{UpperShadow} - \text{LowerShadow}) / (\text{HLSpread} + \epsilon)$
    \item Momentum Acceleration: $\text{ret\_5d} - \text{ret\_21d}$
    \item Volatility-Adjusted Trend: $\text{dist\_sma\_20} / (\text{vol\_21d} + \epsilon)$
\end{itemize}
The primary target is the daily cross-sectional $z$-score of 5-day forward return:
\begin{equation}
    z_{i,t,5} = \frac{R_{i,t,5} - \mu_t(R_{\cdot,t,5})}{\sigma_t(R_{\cdot,t,5})}
\end{equation}

\subsection{Chronological Purged Splits and Test Set Lock}
Splits are strictly chronological with 5-day purge intervals:
\begin{itemize}
    \item Train: 2019-09-26 $\rightarrow$ 2024-03-28 (1,134 days; 2,276,725 rows)
    \item Purge 1: 5 trading days embargo
    \item Validation: 2024-04-08 $\rightarrow$ 2025-06-27 (307 days; 747,545 rows)
    \item Purge 2: 5 trading days embargo
    \item Test: 2025-07-08 $\rightarrow$ 2026-09-25 (308 days; 737,805 rows; LOCKED)
\end{itemize}

\section{Empirical Results}

\subsection{Out-of-Time Forecasting Performance}
Table \ref{tab:forecast} details the forecasting benchmark on the out-of-time test partition.

\begin{table}[H]
\centering
\small
\caption{Out-of-Time Forecasting Performance on Test Partition ($H=5$ Days)}
\label{tab:forecast}
\begin{tabular}{lcccc}
\toprule
\textbf{Model Architecture} & \textbf{Mean Rank IC} & \textbf{IC IR} & \textbf{$t$-statistic} & \textbf{Directional Acc.} \\
\midrule
BASELINE\_MOMENTUM & -0.0234 & -0.189 & -3.28 & 48.97\% \\
OLS\_LINEAR & 0.0015 & 0.015 & 0.26 & 51.06\% \\
RIDGE\_REGRESSION & 0.0015 & 0.015 & 0.26 & 51.06\% \\
RANDOM\_FOREST\_GPU & 0.0005 & 0.005 & 0.08 & 50.32\% \\
XGBOOST\_GBDT\_GPU (30 feats) & 0.0059 & 0.061 & 1.06 & 51.71\% \\
LIGHTGBM\_HUBER (35 feats) & \textbf{0.0085} & 0.057 & 1.48 & 51.49\% \\
MULTI\_MODEL\_STACKING (M4) & \textbf{0.0085} & \textbf{0.075} & 1.51 & 51.42\% \\
\bottomrule
\end{tabular}
\end{table}

\subsection{Top-5 Recommendation Performance}
Table \ref{tab:recommendations} summarizes recommendation performance against benchmarks across 62 out-of-time rebalance dates.

\begin{table}[H]
\centering
\small
\caption{Top-5 Recommendation Performance vs Benchmarks on Test Partition}
\label{tab:recommendations}
\begin{tabular}{lccccc}
\toprule
\textbf{Strategy} & \textbf{Mean Excess} & \textbf{Std Excess} & \textbf{$t$-stat} & \textbf{Wilcoxon $p$} & \textbf{Hit Rate} \\
\midrule
Universe Benchmark & 0.00\% & - & - & - & - \\
Random Top-5 (Monte Carlo) & +0.01\% & 3.43\% & 0.08 & 0.491 & 47.27\% \\
Method A: Prediction-Only & \textbf{+1.66\%} & 12.03\% & 4.83 & \textbf{0.0159} & 44.92\% \\
Method B: Similarity-Only & -0.03\% & 3.91\% & -0.31 & 0.2266 & 48.36\% \\
Method C: Combined Fusion & +0.16\% & \textbf{4.17\%} & 1.33 & 0.9181 & 49.30\% \\
Method C2: Vol-Penalized Fusion & \textbf{+0.59\%} & 9.16\% & \textbf{2.24} & \textbf{0.0248} & \textbf{50.90\%} \\
\bottomrule
\end{tabular}
\end{table}

\section{Discussion and Limitations}
\begin{itemize}
    \item \textbf{Tracking Error Dampening:} Method C dampens excess return volatility by over $65\%$ compared to Method A (from $12.03\%$ to $4.17\%$).
    \item \textbf{Transaction Frictions:} Method B exhibits low turnover ($9.0\%$), while Method C exhibits $83.8\%$ turnover. At round-trip costs $\ge 20$ bps, weekly rebalancing absorbs marginal excess returns.
    \item \textbf{Survivorship Limitation:} Requiring continuous 7-year trading history conditions on survival, representing liquid persistent equities.
\end{itemize}

\section{Conclusion}
This study proves that pure similarity recommendation provides zero alpha, while unconstrained return prediction incurs severe tracking error volatility. Pre-specified rank fusion successfully reconciles this trade-off, delivering positive excess returns with compressed portfolio volatility.

\bibliographystyle{plain}
\bibliography{references}

\end{document}
